% ============================================================

% Plantilla de reporte semanal — Física Computacional (106018C)

% Uso: llenar los campos marcados con \[ ... ] y las secciones.

% Debe caber en UNA página. Si se pasa, es señal de que hay

% que resumir (no de agrandar márgenes ni achicar la fuente).

%

% NOTA PARA QUIEN ARMA EL TALLER (no borrar al reutilizar la plantilla):

% Si el taller de la semana viene dividido en varias partes (p. ej.

% Parte A "métodos/precisión numérica" + Parte B "problema físico" +

% Parte C "verificación"), las secciones 1-3 (contexto físico, método,

% resultado) se anclan SIEMPRE en la parte que tiene un sistema físico

% real y un valor de referencia conocido — no en la parte de métodos

% numéricos por sí sola. Esa parte de métodos no queda fuera del

% reporte: normalmente es ahí donde aparece un obstáculo real (pérdida

% de precisión, criterio de parada, etc.), así que puede alimentar la

% Sección 4 (Obstáculo o decisión de implementación) aunque no sea la

% misma parte que las secciones 1-3. Si el taller no tiene esa

% división (una sola actividad integrada), esta nota no aplica.

% ============================================================

\\documentclass\[11pt]{article}



\\usepackage\[a4paper,margin=2cm]{geometry}

\\usepackage\[spanish]{babel}

\\usepackage\[utf8]{inputenc}

\\usepackage{amsmath,amssymb}

\\usepackage{booktabs}



\\pagestyle{empty}

\\setlength{\\parindent}{0pt}

\\setlength{\\parskip}{4pt}



% ---------- Datos del taller: editar aquí ----------

\\newcommand{\\curso}{Física Computacional (106018C)}

\\newcommand{\\taller}{Taller \[N.º] --- \[1]}

\\newcommand{\\semana}{Semana \[N.º 4]}

\\newcommand{\\lenguaje}{\[Fortran]}

\\newcommand{\\estudiante}{\[Angie Katherine Nazate Cuaspa 202422480]}

\\newcommand{\\fechaentrega}{\[19/09/2025]}

% -----------------------------------------------------



\\begin{document}



\\begin{center}

{\\large\\bfseries \\curso} \\\\\[2pt]

{\\bfseries \\taller} \\\\\[2pt]

\\semana\\ \\quad --- \\quad Lenguaje: \\lenguaje \\\\\[2pt]

\\estudiante\\ \\quad --- \\quad \\fechaentrega

\\end{center}



\\vspace{4pt}

\\hrule

\\vspace{6pt}



% ---------------- 1. Contexto físico ----------------

\\textbf{1. Contexto físico.} % 1--2 frases: qué problema se resuelve y por qué importa.

\[Se estudió el movimiento de un proyectil lanzado con una velocidad inicial y un ángulo determinados, con el objetivo de calcular la altura máxima alcanzada. Este cálculo permite analizar una característica importante del movimiento parabólico mediante herramientas computacionales.]



% ---------------- 2. Método / Algoritmo ----------------

\\textbf{2. Método.} % Fórmula o pseudocódigo clave, no el código completo (eso va en el archivo adjunto).

\[Se utilizó la ecuación de la altura máxima del movimiento parabólico. El programa recibe la velocidad inicial y el ángulo de lanzamiento, convierte el ángulo de grados a radianes y calcula la altura máxima.

]



% Ejemplo: \\int\_a^b f(x)\\,dx \\approx \\frac{h}{2}\\sum\_{i=0}^{n-1}\\left\[f(x\_i)+f(x\_{i+1})\\right]

\[\\begin{equation\*}

h\_{\\max}=\\frac{v\_0^2\\sin^2(\\theta)}{2g}

\\end{equation\*}]



% ---------------- 3. Resultado numérico ----------------



\\textbf{3. Resultado.}% Comparación explícita contra el valor de referencia, con error.

Se realizaron tres casos de prueba y se compararon los resultados obtenidos por el programa con los valores de referencia. \\begin{center} \\begin{tabular}{lccc} \\toprule Caso \& Valor obtenido (m) \& Valor de referencia (m) \& Error relativo \\\\ \\midrule 1: 90° \& 4.9000010 \& 4.9000 \& 0.00002\\% \\\\ 2: 45° \& 10.2040806 \& 10.2041 \& 0.00019\\% \\\\ 3: 30° \& 5.1020477 \& 5.1020 \& 0.00094\\% \\\\ \\bottomrule \\end{tabular} \\end{center} \\textbf{Verificación (PASS/FAIL):} 3 de 3 casos de la celda de verificación pasaron.



% ---------------- 4. Obstáculo / decisión de implementación ----------------

\\textbf{4. Obstáculo o decisión de implementación.} % Puede ser el mismo punto que se publicó en el foro.

\[Un aspecto importante de la implementación fue realizar correctamente

la conversión del ángulo de grados a radianes, ya que las funciones

trigonométricas utilizadas en Fortran trabajan con radianes. Además,

se realizaron casos de prueba para comprobar que los resultados

obtenidos coincidieran con los valores de referencia.]



% ---------------- 5. Conclusión ----------------

\\textbf{5. Conclusión.} % 1--2 frases: qué se aprendió o qué se haría distinto.

\[Se aprendió a implementar una expresión física en Fortran y a verificar numéricamente los resultados mediante casos de prueba. La comparación con los valores de referencia mostró que el programa calcula correctamente la altura máxima para los casos evaluados.]



% ---------------- 6. Uso de herramientas de IA ----------------

\\textbf{6. Uso de herramientas de IA.} % Declarar qué se usó y para qué -- no se penaliza el uso, sí la falta de declaración. Mantener a 1-2 líneas para conservar el límite de una página.

\[Se utilizó ChatGPT como herramienta de apoyo para comprender la implementación del problema en Fortran, revisar errores y organizar la estructura del reporte. La ejecución, compilación y verificación de los resultados se realizaron mediante el código desarrollado y probado en Google Colab.]



\\vspace{4pt}

\\hrule

\\vspace{4pt}

{\\small Código fuente: ver archivo adjunto \\texttt{\[proyectil.f90]}.}



\\end{document}

